\chapter{Conclusions and Further Work}
\label{ch:conclusion}

\section{Summary of Findings}
\label{sec:conclusion-findings}

This project set out to find whether a sound too high-pitched to hear,
combined with strong cryptography, can give a passwordless web login that
needs no typing beyond a username, no aiming a camera, and no second app
to think about, while still standing up to the attacks that make
passwords and their replacements weak. Looking back at the research
questions from Section~\ref{ch:introduction}:

\textbf{RQ1 (does the sound actually work).} The early feasibility test
and the finished build both confirm that an ordinary laptop speaker and an
ordinary phone microphone can reliably exchange a short signal at
18 to 20\,kHz at close range, which is why the project was able to move
forward into full development. What this report cannot yet say with the
same confidence is the exact success rate across the full range of
distances and noise levels planned in Section~\ref{sec:methodology-eval},
because, as stated plainly in Table~\ref{tab:evalresults}, the full
270-attempt test had not been carried out on the real presentation
hardware at the time of writing. This is the single most important piece
of work still to do before the project's performance claim is fully
backed up by evidence, and it is said clearly here rather than covered up
with invented numbers.

\textbf{RQ2 (is the cryptography correct).} This question can be answered
with much more confidence. The 31-check automated test suite
(Section~\ref{sec:testing}), run again for this report, passed completely:
replaying a used code, forging a signature with the wrong key, using a
valid signature against the wrong device, reusing a claim token, and
trying to log in with a removed phone were all correctly blocked, and the
background check correctly stopped one user's phone from ever being asked
to approve someone else's login. Within what the tests actually check, the
cryptography behaves exactly as designed: nothing secret ever needs to
travel through the sound, and a one-time code combined with a
device-specific signature makes the classic replay and forgery attacks
that break weaker passwordless methods (Chapter~\ref{ch:literature})
effectively impossible against this system.

\textbf{RQ3 (how fast is it).} The full process, creating the code,
playing it as sound, decoding it, the background check, user approval,
signing, checking, and confirming the session, was designed around an
8 to 10 second target based on the demo script used during development.
Exact, measured timing figures depend on the same real-world test
described under RQ1, and are left for that completed test rather than
guessed at here.

\textbf{RQ4 (what risks are left).} The clearest weakness is a
\textbf{relay attack}: because the sound channel only shows that a phone
is probably nearby, not that it definitely is, an attacker who can stream
the sound from the victim's laptop to the victim's phone over the
internet, instead of through the air, could trick the system into
approving a login without ever being physically close. Echo does not
currently defend against this. It was treated from the start as a known,
accepted risk rather than a solved problem, because properly solving it
needs very precise timing measurements that a web browser simply does not
allow a developer to make. This is a real weakness, and it is presented
as one, not played down. A second, smaller risk is an unlocked phone that
has been stolen and is physically near the victim's laptop: because Echo
does not currently require a fingerprint or face check on the phone
immediately before signing (a feature that was planned but not finished
in time), anyone holding an unlocked, registered phone near the laptop
could approve a login. Both of these risks, and the reasoning for treating
them as known limitations rather than blockers, are discussed further
below.

\section{Conclusions Against the Objectives}

Measured against the seven objectives set out in
Section~\ref{ch:introduction}, five were fully met: comparing why
passwords fail against the existing passwordless options
(Chapter~\ref{ch:literature}); the sound-based transmission layer with an
audible backup (Section~\ref{sec:architecture}); the full server, covering
sign-up, the one-time code system, signature checking, live updates, and
account recovery (Chapter~\ref{ch:design}); the phone app with its own
signing key (Chapter~\ref{ch:design}); and the automated test suite
(Section~\ref{sec:testing}), which was not only built but was actually run
for this report with a clean 31 out of 31 pass. The seventh objective,
account recovery, was met in a changed form: an email link replaced the
originally planned recovery-code system as the main way to recover an
account, for the reasons given in Section~\ref{sec:design-evolution},
while the original recovery-code system is still present in the code as a
backup option. The sixth objective, testing across distance and noise, is
the one genuinely unfinished item: the tool and the plan for doing it are
built and ready (Section~\ref{sec:methodology-eval}), but the actual
270-attempt data collection still needs to happen on the real
demonstration hardware. Overall, the project delivers a working system
whose cryptography is sound and has been properly tested, but whose
real-world sound reliability is not yet backed by finished measurements,
and that gap is this report's clearest limitation.

\section{Limitations}

Beyond the relay-attack and unlocked-phone risks already discussed, three
more limitations are worth stating plainly. First, once finished, the
real-world test will only describe one specific laptop and phone.
Speaker and microphone quality varies a lot between different devices, so
a result from the development machine should not be assumed to apply to
every laptop and phone without further testing. Second, the phone side
depends on the browser continuing to allow microphone access and running
the app in the background the way the phone's operating system allows;
iPhone's Safari browser is known to behave differently around background
microphone access, and this was not fully solved within the time
available for this project. Third, the system was built and tested as a
single-user, single-phone-per-account demo; it does not yet have a proper
way to register or remove more than one phone per account beyond the
basic version already tested, which would be needed before the design
could be considered ready for real, everyday users.

\section{Recommendations and Further Work}
\label{sec:further-work}

The following are recommended as the next steps, in order of importance:

\begin{enumerate}[label=\arabic*.]
  \item \textbf{Finish the real-world sound test.} Run the full
  270-attempt test described in Section~\ref{sec:methodology-eval} on the
  real presentation hardware, fill in Table~\ref{tab:evalresults} with the
  actual success rates and timing results, and update the RQ1 and RQ3
  conclusions above to match the real numbers before final submission.
  This is the single most important task left.
  \item \textbf{Add a fingerprint or face check before signing.}
  Requiring the phone to check the user's fingerprint or face right
  before it signs a login closes the unlocked-phone risk described above.
  This was already planned as the next priority but did not make it into
  this build, and should be the first thing added if the project
  continues.
  \item \textbf{Look into simple ways to reduce the relay-attack risk.}
  A full, certain fix is out of reach of a web browser, but two simpler
  ideas are worth trying: timing how long the whole process takes as a
  rough distance signal, and comparing a short recording of background
  noise from both the laptop and the phone to check they are likely in the
  same room. Neither would fully solve the problem alone, but either would
  make a relay attack harder than it currently is.
  \item \textbf{Support more than one phone per account, with proper
  management.} The database already allows more than one phone per user;
  building a proper screen for registering and removing phones would move
  the project from a single-user demo toward something closer to a real,
  usable feature.
  \item \textbf{Build a proper phone app.} A dedicated Android or iPhone
  app, using the phone's own secure hardware for storing the signing key,
  would remove the web browser as a weak point for key storage entirely.
  Right now the key cannot be exported, but it still lives inside the
  browser's storage rather than in dedicated secure hardware built for
  exactly this purpose.
\end{enumerate}

\section{Closing Remarks}

The main idea this project tested, that a nearby, high-pitched sound
combined with strong cryptography can be a workable passwordless login
method, holds up well on the cryptography side. Its resistance to replay
attacks, forged signatures, and phishing is not just designed on paper but
has been properly tested and confirmed with a clean, fully passing set of
automated tests. Its real-world sound reliability is meant to be tested
with the same level of care, and the honest state of this report is that
the real-world data needed to fully back that up still needs to be
collected. Stating that gap openly fits the wider goal of this project:
to make every part of a passwordless login, what is secret, what is
public, what expires, and what still needs to be proven, clear and
checkable, instead of hidden behind one company's closed system.
